% ============================================
% 三、阅读程序填空题 Q6-10
% ============================================
\section{阅读程序填空题（二）}

% --- Reading Q6 ---
\begin{frame}[fragile]{阅读程序 第6题}{static 局部变量}
    \begin{exampleblock}{输出结果：\textbf{9}}
    \end{exampleblock}
    \begin{lstlisting}[language=C]
int k = 0, a = 1;
int f3(int c) {
    static int a = 1;
    return (++a) + c;
}
int main() {
    { int a = 3; k += f3(a); }   // 块1
    k += f3(a);                   // 块2
    printf("%d\n", k);
}
\end{lstlisting}
    \pause
    \begin{alertblock}{解析}
        \begin{itemize}
            \item static a 只初始化一次，生命周期贯穿整个程序
            \item 块1：f3(3) → static a从1→2，返回2+3=5，k=5
            \item 块2：f3(a)中a是全局a=1（块1的a已出作用域）\\
            static a从2→3，返回3+1=4，k=5+4=9
        \end{itemize}
        {\color{highlight}关键：} static局部变量\underline{只初始化一次}，值在函数调用间保留！
    \end{alertblock}
\end{frame}

% --- Reading Q7 ---
\begin{frame}[fragile]{阅读程序 第7题}{结构体指针传参}
    \begin{exampleblock}{输出结果：\textbf{x=6, y=5}}
    \end{exampleblock}
    \begin{lstlisting}[language=C]
struct Point { int x; int y; };
void modifyPoint(struct Point *p) {
    int temp = p->x;
    p->x = p->y;                  // x = 12
    p->y = temp;                  // y = 5
    p->x = (p->x > 10) ? (p->x / 2) : (p->x * 2);
}
int main() {
    struct Point myPoint = {5, 12};
    modifyPoint(&myPoint);
    printf("x = %d, y = %d\n", myPoint.x, myPoint.y);
}
\end{lstlisting}
    \pause
    \begin{alertblock}{解析}
        myPoint初始 \{x=5, y=12\}：
        \begin{enumerate}
            \item 交换：x↔y → \{12, 5\}
            \item 条件运算：x=12>10 → x = 12/2 = 6
        \end{enumerate}
        最终：x=6, y=5（通过指针修改了原结构体）
    \end{alertblock}
\end{frame}

% --- Reading Q8-10 ---
\begin{frame}[fragile]{阅读程序 第8-10题（一）}{第8空}
    \begin{exampleblock}{第8空输出：\textbf{z=8}}
    \end{exampleblock}
    \begin{lstlisting}[language=C]
int func(int a, int b) {
    if (b == 0) return 1;
    int result = a;
    for (int i = 1; i < b; i++) result *= a;
    return result;
}
int main() {
    int x = 2, y = 3, z = func(x, y);
    printf("z = %d\n", z);  // 【8】
}
\end{lstlisting}
    \pause
    \begin{alertblock}{解析}
        \texttt{func(2, 3)} 计算 $2^3$：
        \begin{itemize}
            \item result = 2
            \item i=1: result *= 2 → 4
            \item i=2: result *= 2 → 8
        \end{itemize}
        func实现的是幂运算 $a^b$
    \end{alertblock}
\end{frame}

\begin{frame}[fragile]{阅读程序 第8-10题（二）}{第9空}
    \begin{exampleblock}{第9空输出：\textbf{sum=3}}
    \end{exampleblock}
    \begin{lstlisting}[language=C]
int arr[5] = {1, 2, 3, 4, 5};
int sum = 0;
for (int i = 0; i < 5; i++) {
    if (i % 2 == 0) sum += arr[i];
    else sum -= arr[i];
}
printf("sum = %d\n", sum);  // 【9】
\end{lstlisting}
    \pause
    \begin{alertblock}{解析}
        \begin{tabular}{c|c|c|c}
            i & arr[i] & i\%2 & sum变化 \\ \hline
            0 & 1 & 0(偶) & 0+1=1 \\
            1 & 2 & 1(奇) & 1-2=-1 \\
            2 & 3 & 0 & -1+3=2 \\
            3 & 4 & 1 & 2-4=-2 \\
            4 & 5 & 0 & -2+5=3
        \end{tabular}
        偶数下标加，奇数下标减
    \end{alertblock}
\end{frame}

\begin{frame}[fragile]{阅读程序 第8-10题（三）}{第10空：异或交换}
    \begin{exampleblock}{第10空输出：\textbf{a=3, b=5}}
    \end{exampleblock}
    \begin{lstlisting}[language=C]
int a = 5, b = 3;
a = a ^ b; b = a ^ b; a = a ^ b;
printf("a = %d, b = %d\n", a, b);  // 【10】
\end{lstlisting}
    \pause
    \begin{alertblock}{异或交换法（不借助临时变量）}
        二进制视角：
        \begin{itemize}
            \item a=5(101), b=3(011)
            \item a=a\^{}b → a=110(6)
            \item b=a\^{}b → b=101(5)
            \item a=a\^{}b → a=011(3)
        \end{itemize}
        {\color{highlight}技巧：} \texttt{x \^{} y \^{} y = x}（两次异或同一值恢复原值）

        本质：利用 $A \oplus B \oplus B = A$ 的性质
    \end{alertblock}
\end{frame}
